\newpage
\section{假设与符号说明}

\subsection{基本假设}
\begin{enumerate}
    \item 用户充电行为可用需求强度指数表征，忽略个体差异
    \item 充电站建设周期远短于需求变化周期，静态优化可行
    \item 服务半径内用户均会选择最近的充电站
    \item 电网容量充足，建站不受供电能力约束
\end{enumerate}

\subsection{符号说明}
\begin{center}
\begin{tabular}{c l c}
\hline
符号 & 含义 & 单位 \\
\hline
$I_i$ & 需求点$i$的需求强度指数 & - \\
$d_{ij}$ & 需求点$i$到站点$j$的距离 & km \\
$r_{\max}$ & 最大服务半径 & km \\
$x_j$ & 站点$j$是否建设的0-1变量 & - \\
$c_j$ & 站点$j$的建设成本 & 万元 \\
$N_{cov}$ & 覆盖的需求点数量 & 个 \\
$G$ & 基尼系数衡量空间公平性 & - \\
\hline
\end{tabular}
\end{center}
